\section{Large degree bounds for scalar groups with unramified models}
\label{sec:scalar-unramified}

\begin{theorem}[D-BI-14: the scalar-unramified degree bound]
\label{D-BI-14}
Let $K=\mathbf Q(\sqrt{-p})$ for
$p\in\{7,11,19,43,67,163\}$, and suppose a finite-cover
Weil sixfold source of exact character field $K$ has period rank
nine and nonsplit actual polarization.  Suppose an inducing
representation of degree $e$ has faithful image $J$ satisfying
\[
                 J\subset\mathbf C^*N,
                 \qquad N\subset\operatorname{GL}_e(F),
                                                               \tag{BI.45}
\]
where $N$ is finite and the number field $F$ is unramified at $p$.
The field $F$ need not be Galois, and need not contain $K$ or the
inducing character field.  Then
\[
                  e\ge7(p-1).
                                                               \tag{BI.46}
\]
The six bounds, in the displayed order of $p$, are
\[
                   42,\ 70,\ 126,\ 294,\ 462,\ 1134.
\]
If the inducing image is also primitive and tensor indecomposable
in the Guralnick--Saxl convention, then
\[
                  e\ge12(p-1),
                                                               \tag{BI.47}
\]
giving respectively $72,120,216,504,792,1944$.
More generally, if every nonidentity matrix of $J$ has moving rank
at least $q$, the same argument gives
\[
                  e\ge(5q-3)(p-1).                         \tag{BI.48}
\]
The estimate (BI.48) for $q=1$ gives the weaker bound
$2(p-1)$ without invoking D-BI-13's absence of reflections.
\end{theorem}
\begin{proof}
Let $\tau$ denote the inducing representation and
$L=\mathbf Q(\chi_\tau)$.  Induction of characters gives
$K\subset L$.  D-BI-13 says that $J$ has no reflections.
In particular every genuine inertia has moving rank at least two.
There are no degree-one cases; a faithful nontrivial image in
dimension one would consist of reflections.  On the actual
intermediate base of genus $h$, the multiplicity equation is
$D=6-2e(h-1)$.  Genus one has at most three genuine branch points
and at most three marked-base parameters.  For $h\ge2$ the
only numerical possibilities have $h=2,e\le3$ and at most
$3-e$ branch points, hence dimension at most $6-e<9$.
Consequently the base is rational and
\[
 r\ge12,\qquad \sum_j\operatorname{rank}(1-g_j)=2e+6
                                                               \tag{BI.49}
\]
for a generating product-one tuple of genuine inertias.
For the stated $q=1$ variant, one can instead discard $e=1$
by D-MON and use the weaker bound of $D$ genuine branches:
genus one gives at most six parameters, and genus at least two
gives at most five.  Thus (BI.49) still follows without the
reflection-free conclusion of D-BI-13.

We provide the field details for arbitrary $F$.  The extension
$\mathbf Q(\zeta_p)/\mathbf Q$ is totally ramified at $p$:
$\Phi_p(X+1)$ is Eisenstein at $p$, as its constant term is
$p$, all nonleading coefficients are divisible by $p$, and its
leading coefficient is one.  Every nontrivial subfield is therefore
ramified at $p$.  Since $F$ is unramified at $p$,
\[
              F\cap\mathbf Q(\zeta_p)=\mathbf Q,
              \qquad [F(\zeta_p):F]=p-1.                  \tag{BI.50}
\]
For completeness, the degree conclusion can also be seen locally:
base change to an unramified completion of $F$ preserves the
Eisenstein property of $\Phi_p(X+1)$, so it remains irreducible
over that completion and hence over $F$.
This avoids imposing a Galois hypothesis on $F$.

Every scalar in $Z_0=N\cap\mathbf C^*I$ is a root of unity in
$F$ and has order prime to $p$, by (BI.50).  As in D-BI-13,
write $g_j=\lambda_j n_j$ with $n_j\in N$ and a root of unity
$\lambda_j$.  Its $p$-primary part is independent of the choice
modulo $Z_0$, and defines a character $\beta_p$ of $J$.
It is nontrivial: if all generating phases had order prime to
$p$, the field $F(\lambda_1,\ldots,\lambda_r)$ would be
unramified at $p$ and would contain $L$, hence the ramified field
$K$, an impossibility.  Here adjoining roots of unity of order
prime to $p$ is unramified locally, and a compositum of unramified
local extensions is unramified.  The product-one relation supplies
at least two $g_j$ with nontrivial $p$-primary phase.

For either of these two entries, (BI.50) gives
$[F(\lambda_j):F]\ge p-1$.  If its fixed dimension is $b$,
the irreducible polynomial of $\lambda_j^{-1}$ over $F$
occurs at least $b$ times in the characteristic polynomial of
$n_j$.  Thus its moving rank is at least
$e-\lfloor e/(p-1)\rfloor$.  If all other nonidentity entries
have moving rank at least $q$, (BI.49) yields
\[
 2e+6\ge2e-2\left\lfloor\frac{e}{p-1}\right\rfloor+10q,
 \qquad
 \left\lfloor\frac{e}{p-1}\right\rfloor\ge5q-3.
\]
This proves (BI.48), and $q=2$ proves (BI.46).

Finally assume primitivity and tensor indecomposability.  The
already established bound gives $e\ge42$, so the degree-ten
exception in Guralnick--Saxl cannot occur.  If $J$ contained a
matrix of moving rank two, their characteristic-zero classification
\cite[Theorem~7.1 and Remark~(3)]{BI:GS} would make it a scalar
deleted alternating/symmetric representation.  D-BI-08 excludes
that inducing block, with arbitrary larger scalar field.  Thus
every nonidentity element now moves at least three dimensions.
Applying (BI.48) with $q=3$ proves (BI.47).
\end{proof}
